\section{Background and Related Work}
\label{sec:related}

\subsection{A Running Example}
\label{sec:example}

Our test application, \texttt{broken-admin}, has six pages (\texttt{index},
\texttt{login}, \texttt{register}, \texttt{dashboard}, \texttt{reports},
\texttt{admin}); a guest signs in from \texttt{index} and reaches
\texttt{dashboard}, which links to \texttt{reports} and \texttt{admin}, and most
pages link back. The policy has four sentences, each with one planted violation:
(F1/V-01) a guest must never see the admin link, planted on \texttt{dashboard};
(F2/V-02) no page may ship a debug panel, planted on \texttt{reports};
(F3/V-03) a guest must never be able to delete a user, planted as a Delete-user
button on \texttt{admin}; and (R1/V-04) every page must offer a skip-to-content
link, which \texttt{admin}, the only page lacking one, omits. The forbidden defects
are \emph{additional} UI, so a workflow test has no reason to fail, and R1 is an
absence, which an intersection-only oracle cannot express.

\subsection{Related Work}

Table~\ref{tab:related} compares the closest systems we found along the axes of
our literature search (crawling and model inference, LLM/VLM GUI agents,
access-control testing, negation in VLMs, specification and accessibility
oracles) by web search in September and October 2026. It is a positioning aid, not a systematic review, and absence from it is not
evidence that no related system exists. Cells were read from the abstract of each cited source (and from the
body for ATUSA); ``?'' marks what we could not establish.

\begin{table}[tbp]
\caption{Positioning against nearby work. \checkmark{}: stated in the source we
read; $\sim$: partial, or only expressible in code; --: not stated in what we
read; ?: not established (MATE: we could not obtain an abstract naming its checks). AXNav's $\sim$ on REQ reflects behavioural heuristics that flag expected behaviour missing, not required elements. AE: autonomous exploration; POL: declarative policy as
input; NEG: must-not-exist oracle; REQ: required-absent oracle; ROLE: role- or
user-scoped; MM: multimodal perception. \sys{} marks mean implemented and
exercised at fixture level only (AE: scripted crawler; MM: unevaluated, not wired
in; POL: no English compiler).}
\label{tab:related}
\setlength{\tabcolsep}{4pt}
\renewcommand{\arraystretch}{1.02}
\footnotesize
\begin{tabular*}{\textwidth}{@{\extracolsep{\fill}}lcccccc@{}}
\toprule
System & AE & POL & NEG & REQ & ROLE & MM \\
\midrule
Crawljax \cite{mesbah2012crawljax}              & \checkmark & -- & -- & -- & -- & -- \\
ATUSA \cite{mesbah2012invariant}                & \checkmark & $\sim$ & \checkmark & \checkmark & -- & -- \\
WebMate \cite{dallmeier2013webmate}             & \checkmark & -- & -- & -- & -- & -- \\
Kea \cite{xiong2024kea}                         & \checkmark & $\sim$ & $\sim$ & $\sim$ & -- & -- \\
MATE \cite{eler2018mate}                        & \checkmark & -- & ? & ? & -- & -- \\
Sun et al. \cite{sun2011static}                 & -- & -- & $\sim$ & -- & \checkmark & -- \\
DroidAgent \cite{yoon2024droidagent}            & \checkmark & -- & -- & -- & -- & -- \\
VisionDroid/Trident \cite{liu2024visiondroid} & \checkmark & -- & -- & -- & -- & \checkmark \\
AXNav \cite{taeb2024axnav}                      & -- & $\sim$ & -- & $\sim$ & -- & $\sim$ \\
UIGuard \cite{chen2023uiguard}                  & -- & $\sim$ & \checkmark & -- & -- & \checkmark \\
ST-WebAgentBench \cite{levy2026stwebagentbench} & -- & \checkmark & $\sim$ & $\sim$ & -- & -- \\
axe-core \cite{deque2024axecore,deque2024bypass,deque2024marquee} & -- & $\sim$ & \checkmark & \checkmark & -- & -- \\
ZAP \cite{zap2025passive,zap2025users,zap2025spider} & \checkmark & -- & \checkmark & -- & $\sim$ & -- \\
\midrule
\textbf{\sys{}} & $\sim$ & $\sim$ & \checkmark & \checkmark & \checkmark & $\sim$ \\
\bottomrule
\end{tabular*}
\end{table}

No row combines exploration, a declarative role-scoped policy and both polarities
as typed fields; this describes the table, not all prior work. ATUSA supplies both
polarities per application through code. axe-core ships fixed rules of both
polarities (\texttt{bypass} requires a mechanism, \texttt{marquee} forbids an
element) and ZAP ships passive rules and per-user scanning, but neither has
role-scoped, user-declared clauses.

\paragraph{Crawlers and GUI testers.}
GUI Ripping \cite{memon2003ripping}, Crawljax \cite{mesbah2012crawljax} and WebMate
\cite{dallmeier2013webmate} recover state-flow models of running interfaces; on
Android, Sapienz, Stoat, APE, DroidBot and Humanoid
\cite{mao2016sapienz,su2017stoat,gu2019ape,li2017droidbot,li2019humanoid} explore
by search, stochastic models, model refinement or learned input generation, and
Stoat reports crash detection as its headline result. Such models are cyclic by
construction, so a cyclic graph is not our novelty. Our fingerprint is an ordinary
state abstraction \cite{yandrapally2020nearduplicate}, and Doup\'e et al.\
\cite{doupe2012enemy} infer a server-side state machine from state-changing
requests, which our claim rule only echoes for the auditor's own side effects.

\paragraph{Oracles, invariants and deployed web checkers.}
Barr et al.\ \cite{barr2015oracle} survey the oracle problem; Memon et al.\
\cite{memon2000oracles} introduced GUI oracles. ATUSA \cite{mesbah2012invariant}
checks generic and application-specific invariants on every state Crawljax
reaches, written as DOM conditions with XPath and JavaScript preconditions, and
its case studies include invariants that require a structure to exist and
invariants that flag unwanted content. In expressive power our oracle therefore
offers a skilled tester nothing that hand-written ATUSA invariants could not
(Sect.~\ref{sec:rq5} compares with a hand-written script of that kind, our own
re-implementation and not ATUSA). The differences are that polarity is a typed
field, clauses use a vocabulary of roles, names, tags and capabilities in place of
page-specific XPath, clauses carry a role scope, and each finding records the
operation that produced it. Kea \cite{xiong2024kea} checks user-written properties
and MATE \cite{eler2018mate} applies accessibility checks during Android
exploration.

Both polarities are also checked on the web. The axe-core engine
\cite{deque2024axecore} has a \texttt{bypass} rule that passes a page offering an
internal skip link, a heading or a landmark \cite{deque2024bypass}, implementing
WCAG success criterion 2.4.1 \cite{w3c2023wcag241}, and a \texttt{marquee} rule that
forbids a deprecated element \cite{deque2024marquee}. The first is a deployed
required-presence oracle close to our example R1, although it also accepts a heading
or landmark where R1 demands a skip link; both run on whatever page a test driver
supplies. ZAP
has a spider that discovers resources from seed URLs \cite{zap2025spider}, a
passive rule that flags known debug-error fragments in responses
\cite{zap2025passive}, and it can issue requests as a configured user in a
context \cite{zap2025users}. Detecting a missing skip link or a leaked debug
artifact is thus not new; our contribution is treating both polarities uniformly in
one role-scoped, declarative clause language over an explored state graph.

\paragraph{Access-control testing.}
Sun et al.\ \cite{sun2011static} statically build per-role sitemaps, compare them to
find privileged pages, and check whether forced browsing to them succeeds. Li et
al.\ \cite{li2014blackbox}, Noseevich and Petukhov \cite{noseevich2011detecting},
MACE \cite{monshizadeh2014mace}, BACScan \cite{liu2025bacscan} (a feedback-driven
oracle on read and modify operations) and BACFuzz \cite{dharmaadi2025bacfuzz} target
authorization flaws on the server side. \sys{} instead asks whether the interface
shown to a role conforms to a declared policy, including a required-absent half
those tools do not model. Exposure of a link is not by itself an OWASP A01
vulnerability \cite{owasp2021a01}, so we position the work as policy conformance.

\paragraph{LLM and VLM agents, negation, and specifications.}
GPTDroid, DroidAgent and AutoDroid \cite{liu2024gptdroid,yoon2024droidagent,wen2024autodroid}
drive mobile apps with language models; VisionDroid (renamed Trident in a later
arXiv version) \cite{liu2024visiondroid} asks a multimodal model to detect non-crash
functional bugs; AXNav \cite{taeb2024axnav},
AUITestAgent \cite{hu2024auitestagent} and WebTestPilot \cite{teoh2026webtestpilot}
turn natural-language descriptions into interactions and checks; UIGuard
\cite{chen2023uiguard} detects dark patterns, a fixed taxonomy of forbidden UI.
WebArena, Mind2Web, OSWorld and SeeAct
\cite{zhou2024webarena,deng2023mind2web,xie2024osworld,zheng2024seeact} concern web
and computer-use agents, and ST-WebAgentBench
\cite{levy2026stwebagentbench} measures whether \emph{agents} comply with
policies, whereas we use a policy to audit the \emph{application}. In most of
these systems the model delivers the verdict. In \sys{} it is not meant to decide
which elements exist, but it can still influence a verdict through the tags it
adds (Sect.~\ref{sec:perception}). VLMs perform near chance on negated queries
\cite{alhamoud2025negation} and show bag-of-words behaviour
\cite{yuksekgonul2023aro,ma2023crepe}; these studies use natural images, so we
cite them as motivation, and transferring them to UI yes/no questions is an
extrapolation we do not test. nl2spec and NL2LTL
\cite{cosler2023nl2spec,fuggitti2023nl2ltl} translate language into temporal
logic; our planned English-to-clause compiler is not implemented, and our policies
are supplied in structured form. Sridhar et al.\ \cite{sridhar2020firepixel,sridhar2020rule} report fire-detection
work (by title and venue only; we make no claim about their method) from the same
institution; it is listed as adjacent work, not as support for our approach.
